% from ega4-10-fr.tex
\oldpage[IV-3]{111}
\begin{remarks}[10.8.2]
\label{IV.10.8.2-fr}
(i) En tout point fermé $x\in X$, la profondeur rectifiée est égale à la
profondeur.

(ii) Soient $Y$ un sous-préschéma fermé de $X$, $j:Y\to X$ l'injection canonique,
$\mathcal G$ un $\mathcal O_Y$-Module cohérent. On sait (5.7.3, (vi)) que l'on a
$\operatorname{prof}(\mathcal G_x)=\operatorname{prof}(j_*(\mathcal G)_x)$ pour tout
$x\in Y$~; on en déduit que l'on a aussi
$\operatorname{prof}^{*}_{x}(\mathcal G)
=\operatorname{prof}^{*}_{x}(j_*(\mathcal G))$ pour tout $x\in Y$.

(iii) La notion de profondeur rectifiée n'a d'intérêt que lorsqu'elle est de
caractère \emph{local}, c'est-à-dire qu'elle ne change pas lorsqu'on remplace $X$
par un voisinage ouvert arbitraire $U$ de $x$. Cela exige évidemment que $x$ ne
soit pas isolé dans $\overline{\{x\}}$ lorsque $x$ n'est pas fermé et par suite
que $X$ soit un \emph{préschéma de Jacobson} (10.4.5.1)~; le plus souvent, il sera
aussi nécessaire de savoir que $\dim(U)=\dim(X)$ pour tout ouvert dense $U$ dans
$X$, et il faudra donc supposer que $X$ vérifie aussi les conditions $2^\circ$ et
$3^\circ$ de (10.6.1).
\end{remarks}

\begin{lemma}[10.8.3]
\label{IV.10.8.3-fr}
Soient $X$ un préschéma régulier et biéquidimensionnel, $\mathcal F$ un
$\mathcal O_X$-Module cohérent. Alors on a, pour tout $x\in X$,
\begin{equation}
\label{IV.10.8.3.1-fr}
\tag{10.8.3.1}
\operatorname{prof}^{*}_{x}(\mathcal F)
=\dim(X)-\dim.\operatorname{proj}(\mathcal F_x).
\end{equation}
\end{lemma}

En effet, comme $X$ est biéquidimensionnel, on a ($\mathbf 0$, 14.3.5.1)
\[
  \dim(\overline{\{x\}})=\dim(X)-\operatorname{codim}(\overline{\{x\}},X)
  =\dim(X)-\dim(\mathcal O_{X,x})
\]
en vertu de (5.1.2). D'autre part, comme $X$ est régulier, on a par
($\mathbf 0$, 17.3.4)
\[
  \operatorname{prof}(\mathcal F_x)
  =\dim(\mathcal O_{X,x})-\dim.\operatorname{proj}(\mathcal F_x)
\]
d'où le lemme.

\begin{corollary}[10.8.4]
\label{IV.10.8.4-fr}
Sous les hypothèses de (10.8.3), la fonction
$x\mapsto\operatorname{prof}^{*}_{x}(\mathcal F)$ est semi-continue inférieurement.
\end{corollary}

Cela résulte de (10.8.3.1), puisque
$x\mapsto\dim.\operatorname{proj}(\mathcal F_x)$ est semi-continue supérieurement
(6.11.1).

\begin{proposition}[10.8.5]
\label{IV.10.8.5-fr}
Soient $S$ un préschéma localement noethérien, $X$ un préschéma localement de
type fini sur $S$, $\mathcal F$ un $\mathcal O_X$-Module cohérent. On suppose que $S$
vérifie les conditions suivantes~: $1^\circ$ $S$ est un préschéma de Jacobson~;
$2^\circ$ $S$ est régulier~; $3^\circ$ les composantes irréductibles de $S$ sont
équicodimensionnelles. Alors la fonction
$x\mapsto\operatorname{prof}^{*}_{x}(\mathcal F)$ est semi-continue inférieurement
dans $X$~; autrement dit, pour tout entier $n$, l'ensemble $U_n$ des $x\in X$
tels que $\operatorname{prof}^{*}_{x}(\mathcal F)\geq n$ est ouvert.
\end{proposition}

Comme les anneaux locaux $\mathcal O_{S,s}$ de $S$ sont réguliers, ils sont
universellement caténaires (5.6.4)~; autrement dit, $S$ vérifie les conditions
$1^\circ$, $2^\circ$ et $3^\circ$ de (10.6.1), donc il en est de même de $X$
(10.6.1, (i)). La notion de profondeur rectifiée étant alors de caractère local
(10.8.2, (iii)), on peut se borner au cas où $S=\operatorname{Spec}(A)$ et
$X=\operatorname{Spec}(B)$ sont affines, $A$ étant un anneau régulier et $B$ une
$A$-algèbre de type fini, donc quotient d'un anneau de polynômes
$C=A[T_1,\ldots,T_n]$, et ce dernier est régulier ($\mathbf 0$, 17.3.7). On peut
donc supposer que $X$ est un sous-préschéma fermé d'un préschéma régulier $Y$
vérifiant également les conditions $1^\circ$, $2^\circ$ et $3^\circ$ de (10.6.1)~;
compte tenu de la remarque (10.8.2, (ii)), on est ainsi ramené au cas où $X$ est
en outre régulier et noethérien. Mais comme les anneaux locaux de $X$ sont alors
intègres, les composantes irréductibles de $X$ sont

